% Direct-authority diplomatic French for printed p.196.
\clearpage
\SourcePage{196}

\LayerTarget{tohoku:result:abelian-g-sheaves-additive}
En particulier, les $G$-faisceaux abéliens forment une catégorie additive,
dont nous allons indiquer les propriétés.

\LayerTarget{tohoku:definition:g-o-module}
Plus généralement, soit $O$ un $G$-faisceau d'anneaux avec unité, considérons
les faisceaux $A$ sur $X$ qui sont à la fois des $G$-faisceaux abéliens et des
$O$-mcdules, les opérations de $O$ sur $A$ étant compatibles avec les
opérations de $O$ (i.e. l'homomorphisme naturel de faisceaux d'ensembles
$O\times A\longrightarrow A$ étant un $G$-homomorphisme) : un tel faisceau
sera appelé un \textit{$G$-$O$-Module}. Appelant
\textit{$G$-$O$-homomorphisme} un homomorphisme de faisceaux qui est un
$O$-homomorphisme et un $G$-homomorphisme, on voit que la somme ou le composé
de deux $G$-$O$-homomorphismes est un $G$-$O$-homomorphisme.

\LayerTarget{tohoku:definition:category-g-o-modules}
Donc les $G$-$O$-Modules forment une \textit{catégorie additive} notée
$C^{O(G)}$ ; si $O$ est le faisceau constant des entiers (avec les opérations
triviales de $G$) on obtient de nouveau la catégorie des $G$-faisceaux
abéliens, notée $C^{X(G)}$.

\LayerTarget{tohoku:comparison:trivial-action-g-o-modules}
Si $G$ opère trivialement sur $X$, la catégorie $C^{O(G)}$ s'interprète comme
la catégorie des $O'$-Modules, où $O'$ est un faisceau d'anneaux convenable,
si par exemple $G$ opère trivialement sur $O$, ou aura
$O'=O\otimes_{\mathbf Z}\mathbf Z(G)$, (où $\mathbf Z(G)$ désigne l'algèbre
de $G$ par rapport à l'anneau $\mathbf Z$ des entiers, ou plutôt le faisceau
constant qu'il définit). Les résultats de 3.1 restent valables à de petites
modifications près :

\LayerTarget{tohoku:proposition:5.1.1}
\noindent\textsc{Proposition 5.1.1.} \textit{Soit $O$ un $G$-faisceau
d'anneaux sur $X$. Alors la catégorie additive $C^{O(G)}$ des
$G$-$O$-Modules est une catégorie abélienne, satisfaisant les axiomes AB 5)
et AB 3*) de 1.5, et admet un générateur.}\textsuperscript{9)}

\LayerTarget{tohoku:proof:5.1.1:p1}
La vérification est triviale, sauf l'existence du générateur, dont nous allons
donner une construction maintenant, généralisant celle du 3.1.

\LayerTarget{tohoku:construction:generator-l-u}
Pour tout ouvert $U$ de $X$, soit $L(U)$ le $O$-module somme directe des
$O$-modules $O_{g\cdot U}$ pour $g\in G$. $L(U)$ peut être considéré comme
$G$-faisceau de façon évidente, de sorte que $L(U)$ devient même un
$G$-$O$-Module. Si $A$ est un $G$-$O$-Module quelconque, un
$G$-$O$-homomorphisme de $L(U)$ dans $A$ est connu quand on connaît sa
restriction à $O_U$, qui est un $O$-homomorphisme de $O_U$ dans $A$, qu'on
peut d'ailleurs se donner à l'avance arbitrairement. La donnée d'un tel
homomorphisme équivaut aussi à la donnée d'une section de $A$ sur $U$.

\LayerTarget{tohoku:result:generator-family-g-o-modules}
Si $B$ est un sous-$G$-$O$-Module de $A$ distinct de $A$, on peut trouver sur
un ouvert convenable $U$ une section de $A$ qui n'est pas une section de $A$.
Par suite, la famille de $L(U)$ est une famille de générateurs de $C^{O(G)}$,
et leur somme directe (pour $U$ parcourant tous les ouverts de $X$) est donc
un générateur.

\LayerTarget{tohoku:conclusion:injective-g-o-embedding:p1}
De la proposition 4.1 résulte en particulier que tout $G$-$O$-faisceau est
isomorphe à un sous-faisceau d'un $G$-$O$-Module injectif. Il importe pour la

\LayerTarget{tohoku:footnote:9}
\begin{quote}
\footnotesize
9) Cette proposition, ainsi que divers résultats de la suite, peuvent
s'énoncer aussi dans le contexte plus général qui suit. On se donne une
catégorie $C$, un groupe $G$, et une ``représentation de $G$ par des foncteurs
dans $C$'' : à tout $g\in G$ est associé un foncteur
$F_g:C\longrightarrow C$, de telle façon qu'on ait $F_e=$ foncteur identique,
$F_gF_{g'}=F_{gg'}$ (du moins à des isomorphismes fonctoriels donnés près,
satisfaisant à certaines conditions de cohérence que nous n'expliciterons
pas). Alors on peut construire la catégorie $C^G$ des ``objets de $C$ avec
groupe d'opérateurs $G$'', tout comme $C^{X(G)}$ (pour un espace à groupe
d'opérateurs $X(G)$) peut se construire à partir de $C^X$. (Comparer aussi
avec le cas particulier, exemple 1.7. f). Beaucoup de propriétés vraies pour
$C$ se conservent alors pas passage à $C^G$.
\end{quote}

\TohokuContinuation{197}
